\section{Deep Learning}
\setcounter{aufgabe}{0}

\subsection{Deep Neural Networks}

% ===================================================================
% Aufgaben zu Deep Learning (Keras)
% Themen: Tensoren, Dense-Layer-Parameter, ReLU,
%         Binary/Categorical Crossentropy, Softmax,
%         Metriken bei unausgewogenen Klassen,
%         Optimizer (Momentum, RMSprop, Adam)
% ===================================================================

% -------------------------------------------------------------------
% 1) Tensoren: Form und Anzahl der Eintraege
% -------------------------------------------------------------------
\aufgabe{
Ein Datensatz besteht aus $60\,000$ RGB-Bildern der Größe
$32 \times 32$ Pixel mit drei Farbkanälen.

\begin{enumerate}[label=\alph*)]
  \item Geben Sie den Rang und die Form (shape) des Tensors an, der den
        \textbf{gesamten} Datensatz enthält.
  \item Wie viele einzelne Zahlenwerte enthält dieser Tensor?
  \item Für das Training mit einem Fully-connected-Netz wird jedes Bild zu
        einem Vektor \glqq plattgedrückt\grqq{} (geflattet). Welche Form hat
        der Eingabetensor eines einzelnen Bildes dann?
\end{enumerate}
}

\loesung{
\begin{enumerate}[label=\alph*)]
  \item Der Datensatz ist ein \textbf{Rang-4-Tensor} mit der Form
  \[
    (60\,000,\; 32,\; 32,\; 3) .
  \]
  Die Achsen sind Bildindex, Höhe, Breite und Kanäle.

  \item Anzahl der Einträge:
  \[
    60\,000 \cdot 32 \cdot 32 \cdot 3 = 184\,320\,000 .
  \]

  \item Ein einzelnes Bild wird zu einem Vektor der Länge
  $32 \cdot 32 \cdot 3 = 3072$, also Form $(3072,)$ (Rang 1).
\end{enumerate}
}

% -------------------------------------------------------------------
% 2) Dense-Layer: Anzahl der Parameter
% -------------------------------------------------------------------
\aufgabe{
Ein Fully-connected-Netz (nur Dense-Layer, jeweils mit Bias) hat folgende
Struktur:
\[
  784 \;\rightarrow\; 300 \;\rightarrow\; 100 \;\rightarrow\; 10 .
\]
Für eine Dense-Schicht mit $n_{\text{in}}$ Eingängen und
$n_{\text{out}}$ Neuronen gilt
\[
  \#\text{Gewichte} = n_{\text{in}}\cdot n_{\text{out}}, \qquad
  \#\text{Bias} = n_{\text{out}} .
\]

\begin{enumerate}[label=\alph*)]
  \item Berechnen Sie die Parameterzahl jeder einzelnen Schicht.
  \item Geben Sie die Gesamtzahl der trainierbaren Parameter an.
  \item Wie viele Parameter würden wegfallen weg, wenn man in allen Schichten den
        Bias deaktiviert (\texttt{use\_bias=False})?
\end{enumerate}
}

\loesung{
\begin{enumerate}[label=\alph*)]
  \item Pro Schicht Gewichte plus Bias:
  \[
    784 \rightarrow 300: \quad 784\cdot 300 + 300 = 235\,200 + 300 = 235\,500 ,
  \]
  \[
    300 \rightarrow 100: \quad 300\cdot 100 + 100 = 30\,000 + 100 = 30\,100 ,
  \]
  \[
    100 \rightarrow 10: \quad 100\cdot 10 + 10 = 1\,000 + 10 = 1\,010 .
  \]

  \item Gesamt:
  \[
    235\,500 + 30\,100 + 1\,010 = \textbf{266\,610} .
  \]

  \item Ohne Bias entfielen $300 + 100 + 10 = \textbf{410}$ Parameter.
\end{enumerate}
}

% -------------------------------------------------------------------
% 3) Dense-Schicht: Vorwaertsrechnung mit ReLU
% -------------------------------------------------------------------
\aufgabe{
Eine Dense-Schicht berechnet
$\text{output} = a(\text{input}\cdot W + b)$. Gegeben sind
\[
  \text{input} = \begin{pmatrix} 1 & 2 \end{pmatrix},
  \qquad
  W = \begin{pmatrix} 1 & -2 & 0 \\ 0 & 1 & 3 \end{pmatrix},
  \qquad
  b = \begin{pmatrix} 1 & 0 & -1 \end{pmatrix},
\]
und $a = \text{ReLU}$ mit $\text{ReLU}(x) = \max(0, x)$.

\begin{enumerate}[label=\alph*)]
  \item Berechnen Sie den Vektor $\text{input}\cdot W + b$.
  \item Wenden Sie die ReLU-Funktion komponentenweise an.
\end{enumerate}
}

\loesung{
\begin{enumerate}[label=\alph*)]
  \item Matrix-Vektor-Produkt (Zeilenvektor mal Matrix):
  \[
    \text{input}\cdot W =
    \begin{pmatrix} 1\cdot 1 + 2\cdot 0 & 1\cdot(-2) + 2\cdot 1 & 1\cdot 0 + 2\cdot 3 \end{pmatrix}
    = \begin{pmatrix} 1 & 0 & 6 \end{pmatrix} .
  \]
  Addition des Bias:
  \[
    \text{input}\cdot W + b =
    \begin{pmatrix} 1+1 & 0+0 & 6-1 \end{pmatrix}
    = \begin{pmatrix} 2 & 0 & 5 \end{pmatrix} .
  \]

  \item ReLU komponentenweise:
  \[
    \text{ReLU}\begin{pmatrix} 2 & 0 & 5 \end{pmatrix}
    = \begin{pmatrix} 2 & 0 & 5 \end{pmatrix} .
  \]
  Alle Werte sind hier bereits $\geq 0$, ReLU verändert sie nicht.
\end{enumerate}
}

% -------------------------------------------------------------------
% 4) Binary Crossentropy
% -------------------------------------------------------------------
\aufgabe{
Für die binäre Klassifikation ist der Loss eines einzelnen Datensatzes
\[
  L(y, \hat{y}) = -\bigl[\, y\log(\hat{y}) + (1-y)\log(1-\hat{y}) \,\bigr]
\]
(natürlicher Logarithmus). Berechnen Sie den Loss für folgende Fälle:

\begin{enumerate}[label=\alph*)]
  \item $y = 1$, $\hat{y} = 0{,}9$ (sicher und richtig).
  \item $y = 1$, $\hat{y} = 0{,}1$ (sicher und falsch).
  \item Berechnen Sie den mittleren Loss über die drei Datensätze
        $(y, \hat{y}) \in \{(1;\,0{,}8),\,(0;\,0{,}2),\,(0;\,0{,}7)\}$.
\end{enumerate}
Hinweis: $\log(0{,}9) \approx -0{,}105$, $\log(0{,}1) \approx -2{,}303$,
$\log(0{,}8) \approx -0{,}223$, $\log(0{,}3) \approx -1{,}204$.
}

\loesung{
\begin{enumerate}[label=\alph*)]
  \item Für $y = 1$ bleibt nur $-\log(\hat{y})$:
  \[
    L = -\log(0{,}9) \approx 0{,}105 .
  \]

  \item Ebenfalls $y = 1$, aber $\hat{y} = 0{,}1$:
  \[
    L = -\log(0{,}1) \approx 2{,}303 .
  \]
  Die sichere Fehlklassifikation wird also rund $22$-mal stärker bestraft.

  \item Einzelverluste:
  \[
    (1;\,0{,}8): \; -\log(0{,}8) \approx 0{,}223,
  \]
  \[
    (0;\,0{,}2): \; -\log(1-0{,}2) = -\log(0{,}8) \approx 0{,}223,
  \]
  \[
    (0;\,0{,}7): \; -\log(1-0{,}7) = -\log(0{,}3) \approx 1{,}204 .
  \]
  Mittelwert:
  \[
    L = \frac{0{,}223 + 0{,}223 + 1{,}204}{3} \approx \frac{1{,}650}{3}
      \approx 0{,}550 .
  \]
\end{enumerate}
}

% -------------------------------------------------------------------
% 5) Softmax
% -------------------------------------------------------------------
\aufgabe{
Die Softmax-Funktion wandelt Logits in eine Wahrscheinlichkeitsverteilung um:
\[
  \text{softmax}(z)_k = \frac{e^{z_k}}{\sum_{j=1}^{K} e^{z_j}} .
\]
Gegeben seien die Logits $z = (1;\; 0;\; -1)$.

\begin{enumerate}[label=\alph*)]
  \item Berechnen Sie die Softmax-Wahrscheinlichkeiten $\hat{y}_0$,
        $\hat{y}_1$, $\hat{y}_2$.
  \item Zeigen Sie, dass $\sum_k \hat{y}_k = 1$ gilt.
  \item Welche Klasse sagt das Netz vorher ($\arg\max_k \hat{y}_k$)?
\end{enumerate}
Hinweis: $e^{1} \approx 2{,}718$, $e^{0} = 1$, $e^{-1} \approx 0{,}368$.
}

\loesung{
\begin{enumerate}[label=\alph*)]
  \item Summe der exponenzierten Logits:
  \[
    S = 2{,}718 + 1 + 0{,}368 = 4{,}086 .
  \]
  \[
    \hat{y}_0 = \frac{2{,}718}{4{,}086} \approx 0{,}665, \quad
    \hat{y}_1 = \frac{1}{4{,}086} \approx 0{,}245, \quad
    \hat{y}_2 = \frac{0{,}368}{4{,}086} \approx 0{,}090 .
  \]

  \item Summe:
  \[
    0{,}665 + 0{,}245 + 0{,}090 = 1 .
  \]
  Die Softmax-Ausgabe ist also eine gültige Wahrscheinlichkeitsverteilung.

  \item Die höchste Wahrscheinlichkeit liegt bei $\hat{y}_0$, das Netz sagt
  \textbf{Klasse 0} vorher.
\end{enumerate}
}

% -------------------------------------------------------------------
% 6) Categorical Crossentropy mit Softmax
% -------------------------------------------------------------------
\aufgabe{
Ein Netz mit drei Klassen gibt für einen Datensatz die
Softmax-Wahrscheinlichkeiten
\[
  \hat{\mathbf{y}} = (0{,}7;\; 0{,}2;\; 0{,}1)
\]
aus. Das wahre Label ist One-Hot-codiert. Der Loss ist
\[
  L(\mathbf{y}, \hat{\mathbf{y}}) = -\,\mathbf{y}\cdot\log(\hat{\mathbf{y}}) .
\]

\begin{enumerate}[label=\alph*)]
  \item Berechnen Sie $L$ für das wahre Label $\mathbf{y} = (1;\,0;\,0)$.
  \item Berechnen Sie $L$ für das wahre Label $\mathbf{y} = (0;\,0;\,1)$.
  \item Erklären Sie, warum bei One-Hot-Labels stets nur ein Summand übrig
        bleibt.
\end{enumerate}
Hinweis: $\log(0{,}7) \approx -0{,}357$, $\log(0{,}1) \approx -2{,}303$.
}

\loesung{
\begin{enumerate}[label=\alph*)]
  \item Nur die erste Komponente von $\mathbf{y}$ ist $1$:
  \[
    L = -\log(\hat{y}_0) = -\log(0{,}7) \approx 0{,}357 .
  \]

  \item Hier ist die dritte Komponente $1$:
  \[
    L = -\log(\hat{y}_2) = -\log(0{,}1) \approx 2{,}303 .
  \]
  Dieselbe Vorhersage ist also ein guter Treffer, wenn Klasse $0$ wahr ist,
  aber ein schlechter, wenn Klasse $2$ wahr ist.

  \item Da $\mathbf{y}$ genau eine $1$ (an der wahren Klasse) und sonst
  Nullen enthält, verschwinden im Skalarprodukt
  $\mathbf{y}\cdot\log(\hat{\mathbf{y}})$ alle Terme außer dem der wahren
  Klasse. Es bleibt $L = -\log(\hat{y}_{\text{wahr}})$.
\end{enumerate}
}

% -------------------------------------------------------------------
% 7) Metriken bei unausgewogenen Klassen
% -------------------------------------------------------------------
\aufgabe{
Ein binärer Klassifikator wurde auf $10\,000$ Testbildern ausgewertet.
Davon sind $1\,000$ T-Shirts (Klasse positiv) und $9\,000$ Nicht-T-Shirts.
Die Confusion Matrix ergab:
\[
  \begin{array}{c|c|c}
     & \text{vorhergesagt: kein T-Shirt} & \text{vorhergesagt: T-Shirt}\\\hline
    \text{wahr: kein T-Shirt} & 8\,820 & 180\\
    \text{wahr: T-Shirt}      & 320    & 680
  \end{array}
\]

\begin{enumerate}[label=\alph*)]
  \item Berechnen Sie die Accuracy (Anteil korrekt klassifizierter Bilder).
  \item Berechnen Sie die Fehlerrate je Klasse und daraus die
        \textbf{Balanced Error Rate}
        $\text{BER} = 1 - \tfrac{1}{2}(\text{Sensitivität} + \text{Spezifität})$.
  \item Welche Accuracy erreicht das naive Modell, das \textbf{immer}
        \glqq kein T-Shirt\grqq{} vorhersagt? Vergleichen Sie mit a).
\end{enumerate}
}

\loesung{
\begin{enumerate}[label=\alph*)]
  \item Korrekt sind die Diagonaleinträge:
  \[
    \text{Accuracy} = \frac{8\,820 + 680}{10\,000} = \frac{9\,500}{10\,000}
    = 0{,}95 = \textbf{95\,\%} .
  \]

  \item Klassenweise Trefferraten:
  \[
    \text{Sensitivität (T-Shirt)} = \frac{680}{1\,000} = 0{,}68,
  \]
  \[
    \text{Spezifität (kein T-Shirt)} = \frac{8\,820}{9\,000} = 0{,}98 .
  \]
  Balanced Error Rate:
  \[
    \text{BER} = 1 - \tfrac{1}{2}(0{,}68 + 0{,}98)
    = 1 - 0{,}83 = \textbf{0{,}17} .
  \]

  \item Das naive Modell klassifiziert alle $9\,000$ Nicht-T-Shirts richtig
  und alle $1\,000$ T-Shirts falsch:
  \[
    \text{Accuracy}_{\text{naiv}} = \frac{9\,000}{10\,000} = 0{,}90 .
  \]
  Die $95\,\%$ liegen also nur knapp über der Baseline von $90\,\%$. Bei
  unausgewogenen Klassen ist die Accuracy irreführend, die BER (naiv $0{,}5$)
  entlarvt den Unterschied deutlicher.
\end{enumerate}
}

% -------------------------------------------------------------------
% 8) Batches und Epochen
% -------------------------------------------------------------------
\aufgabe{
Ein Netz wird mit \texttt{model.fit(..., epochs=30, batch\_size=600,
validation\_split=0.25)} trainiert. Der Trainingsdatensatz umfasst vor dem
Split $60\,000$ Bilder.

\begin{enumerate}[label=\alph*)]
  \item Wie viele Bilder werden nach dem Validierungs-Split tatsächlich zum
        Trainieren verwendet?
  \item Wie viele Batches werden pro Epoche verarbeitet?
  \item Wie viele Gewichts-Updates finden über das gesamte Training statt?
\end{enumerate}
}

\loesung{
\begin{enumerate}[label=\alph*)]
  \item Bei $\texttt{validation\_split} = 0{,}25$ werden $25\,\%$
  abgezweigt:
  \[
    60\,000\cdot(1 - 0{,}25) = 45\,000 \text{ Trainingsbilder} .
  \]

  \item Batches pro Epoche:
  \[
    \frac{45\,000}{600} = 75 .
  \]

  \item Ein Update pro Batch, also
  \[
    75\cdot 30 = \textbf{2\,250} \text{ Gewichts-Updates} .
  \]
\end{enumerate}
}

% -------------------------------------------------------------------
% 9) Momentum-Optimizer
% -------------------------------------------------------------------
\aufgabe{
Der Momentum-Optimizer aktualisiert ein Gewicht nach
\[
  v_t = \beta\, v_{t-1} + (1-\beta)\, g_t, \qquad
  w_{t+1} = w_t - \eta\, v_t .
\]
Gegeben sind $\beta = 0{,}9$, $\eta = 0{,}1$, $w_0 = 1$, $v_0 = 0$ und die
Gradienten $g_1 = 2$, $g_2 = 2$, $g_3 = 2$ (konstante Richtung).

\begin{enumerate}[label=\alph*)]
  \item Berechnen Sie $v_1$, $v_2$, $v_3$ und die zugehörigen Gewichte.
  \item Vergleichen Sie den Gesamtschritt $w_0 - w_3$ mit dem Ergebnis von
        reinem SGD ($w_{t+1} = w_t - \eta\, g_t$) über drei Schritte.
\end{enumerate}
}

\loesung{
\begin{enumerate}[label=\alph*)]
  \item Mit $\beta = 0{,}9$, also $(1-\beta) = 0{,}1$:
  \[
    v_1 = 0{,}9\cdot 0 + 0{,}1\cdot 2 = 0{,}2,
    \qquad w_1 = 1 - 0{,}1\cdot 0{,}2 = 0{,}98 ,
  \]
  \[
    v_2 = 0{,}9\cdot 0{,}2 + 0{,}1\cdot 2 = 0{,}38,
    \qquad w_2 = 0{,}98 - 0{,}1\cdot 0{,}38 = 0{,}942 ,
  \]
  \[
    v_3 = 0{,}9\cdot 0{,}38 + 0{,}1\cdot 2 = 0{,}542,
    \qquad w_3 = 0{,}942 - 0{,}1\cdot 0{,}542 = 0{,}8878 .
  \]

  \item Momentum-Gesamtschritt: $w_0 - w_3 = 1 - 0{,}8878 = 0{,}1122$.
  Reines SGD:
  \[
    w_3 = 1 - 3\cdot(0{,}1\cdot 2) = 1 - 0{,}6 = 0{,}4,
    \qquad w_0 - w_3 = 0{,}6 .
  \]
  Momentum baut seinen Schritt erst allmählich auf (der gleitende
  Mittelwert $v_t$ nähert sich langsam dem Gradienten $2$), während SGD
  sofort den vollen Gradienten nutzt. Bei \textbf{konstanter} Richtung würde
  $v_t$ über viele Schritte gegen $2$ konvergieren und Momentum dann ebenso
  schnell wie SGD werden.
\end{enumerate}
}

% -------------------------------------------------------------------
% 10) Adam mit Bias-Korrektur
% -------------------------------------------------------------------
\aufgabe{
Der Adam-Optimizer nutzt
\[
  m_t = \beta_1 m_{t-1} + (1-\beta_1)\, g_t, \qquad
  v_t = \beta_2 v_{t-1} + (1-\beta_2)\, g_t^2,
\]
\[
  \hat{m}_t = \frac{m_t}{1-\beta_1^{\,t}}, \qquad
  \hat{v}_t = \frac{v_t}{1-\beta_2^{\,t}}, \qquad
  w_{t+1} = w_t - \frac{\eta}{\sqrt{\hat{v}_t} + \varepsilon}\,\hat{m}_t .
\]
Gegeben sind $\beta_1 = 0{,}9$, $\beta_2 = 0{,}999$, $\eta = 0{,}1$,
$\varepsilon = 0$, $m_0 = v_0 = 0$, $w_0 = 1$ und der erste Gradient
$g_1 = 0{,}5$.

\begin{enumerate}[label=\alph*)]
  \item Berechnen Sie $m_1$ und $v_1$.
  \item Berechnen Sie die bias-korrigierten Momente $\hat{m}_1$ und
        $\hat{v}_1$.
  \item Berechnen Sie das aktualisierte Gewicht $w_1$.
\end{enumerate}
}

\loesung{
\begin{enumerate}[label=\alph*)]
  \item Erste Momente:
  \[
    m_1 = 0{,}9\cdot 0 + 0{,}1\cdot 0{,}5 = 0{,}05 ,
  \]
  \[
    v_1 = 0{,}999\cdot 0 + 0{,}001\cdot 0{,}5^2 = 0{,}001\cdot 0{,}25
        = 0{,}00025 .
  \]

  \item Bias-Korrektur mit $t = 1$, also $1-\beta_1^{1} = 0{,}1$ und
  $1-\beta_2^{1} = 0{,}001$:
  \[
    \hat{m}_1 = \frac{0{,}05}{0{,}1} = 0{,}5,
    \qquad
    \hat{v}_1 = \frac{0{,}00025}{0{,}001} = 0{,}25 .
  \]
  Die Korrektur hebt die zum Nullpunkt hin verzerrten Momente wieder auf
  ihre eigentliche Größenordnung.

  \item Update:
  \[
    w_1 = 1 - \frac{0{,}1}{\sqrt{0{,}25}}\cdot 0{,}5
        = 1 - \frac{0{,}1}{0{,}5}\cdot 0{,}5
        = 1 - 0{,}1 = 0{,}9 .
  \]
  Der erste Adam-Schritt hat trotz kleinem Gradienten die volle Größe
  $\eta$, da $\hat{m}_1$ und $\sqrt{\hat{v}_1}$ nach der Korrektur genau den
  Gradienten $g_1$ widerspiegeln.
\end{enumerate}
}

% ===================================================================
% Aufgaben zu tiefen Netzen (Deep Learning)
% Themen: Parameterzahl tiefer Netze, Overfitting-Diagnose ueber
%         Train/Test-Abstand, Rechenaufwand der Matrixoperationen,
%         Parallelisierung auf der GPU
% ===================================================================

% -------------------------------------------------------------------
% 1) Parameterzahl eines tiefen Netzes
% -------------------------------------------------------------------
\aufgabe{
Ein tiefes Fully-connected-Netz für Fashion-MNIST hat die Struktur
\[
  784 \;\rightarrow\; 512 \;\rightarrow\; 256 \;\rightarrow\; 128
       \;\rightarrow\; 10 ,
\]
jede Dense-Schicht mit Bias. Für eine Schicht mit $n_{\text{in}}$ Eingängen
und $n_{\text{out}}$ Neuronen gilt
\[
  \#\text{Parameter} = n_{\text{in}}\cdot n_{\text{out}} + n_{\text{out}} .
\]

\begin{enumerate}[label=\alph*)]
  \item Berechnen Sie die Parameterzahl jeder Schicht.
  \item Geben Sie die Gesamtzahl der trainierbaren Parameter an.
  \item Vergleichen Sie mit dem flachen Netz
        $784 \rightarrow 300 \rightarrow 10$. Um welchen Faktor hat das
        tiefe Netz mehr Parameter?
\end{enumerate}
}

\loesung{
\begin{enumerate}[label=\alph*)]
  \item Schichtweise:
  \[
    784 \rightarrow 512: \quad 784\cdot 512 + 512 = 401\,408 + 512 = 401\,920 ,
  \]
  \[
    512 \rightarrow 256: \quad 512\cdot 256 + 256 = 131\,072 + 256 = 131\,328 ,
  \]
  \[
    256 \rightarrow 128: \quad 256\cdot 128 + 128 = 32\,768 + 128 = 32\,896 ,
  \]
  \[
    128 \rightarrow 10: \quad 128\cdot 10 + 10 = 1\,280 + 10 = 1\,290 .
  \]

  \item Gesamt:
  \[
    401\,920 + 131\,328 + 32\,896 + 1\,290 = \textbf{567\,434} .
  \]

  \item Flaches Netz:
  \[
    (784\cdot 300 + 300) + (300\cdot 10 + 10)
    = 235\,500 + 3\,010 = 238\,510 .
  \]
  Faktor:
  \[
    \frac{567\,434}{238\,510} \approx 2{,}38 .
  \]
  Das tiefe Netz hat also gut das $2{,}4$-fache an Parametern. Mehr
  Parameter bedeuten nicht automatisch bessere Ergebnisse, wie das
  Overfitting-Beispiel zeigt.
\end{enumerate}
}

% -------------------------------------------------------------------
% 2) Overfitting-Diagnose ueber den Train/Test-Abstand
% -------------------------------------------------------------------
\aufgabe{
Für drei Netze wurden Trainings- und Testgenauigkeit gemessen:
\[
  \begin{array}{c|c|c}
    \text{Netz} & \text{Trainings-Acc.} & \text{Test-Acc.}\\\hline
    A \;(\text{tief}) & 0{,}982 & 0{,}838\\
    B \;(\text{flach}) & 0{,}962 & 0{,}897\\
    C \;(\text{zu klein}) & 0{,}612 & 0{,}605
  \end{array}
\]

\begin{enumerate}[label=\alph*)]
  \item Berechnen Sie für jedes Netz den Abstand
        $\Delta = \text{Train} - \text{Test}$ in Prozentpunkten.
  \item Ordnen Sie jedem Netz die passende Diagnose zu: \glqq Overfitting\grqq{},
        \glqq gute Verallgemeinerung\grqq{}, \glqq mehr Daten nötig\grqq{}.
  \item Welches Netz würden Sie einsetzen und warum?
\end{enumerate}
}

\loesung{
\begin{enumerate}[label=\alph*)]
  \item Abstände:
  \[
    \Delta_A = 0{,}982 - 0{,}838 = 0{,}144 = 14{,}4 \text{ Prozentpunkte} ,
  \]
  \[
    \Delta_B = 0{,}962 - 0{,}897 = 0{,}065 = 6{,}5 \text{ Prozentpunkte} ,
  \]
  \[
    \Delta_C = 0{,}612 - 0{,}605 = 0{,}007 = 0{,}7 \text{ Prozentpunkte} .
  \]

  \item Diagnosen:
  \[
    A: \text{großer Abstand} \Rightarrow \textbf{Overfitting} ,
  \]
  \[
    B: \text{kleiner Abstand, hohe Werte} \Rightarrow
       \textbf{gute Verallgemeinerung} ,
  \]
  \[
    C: \text{kleiner Abstand, beide niedrig} \Rightarrow
       \textbf{mehr Daten nötig} .
  \]

  \item Netz $B$: Es hat die höchste Testgenauigkeit ($89{,}7\,\%$) bei
  kleinem Abstand zur Trainingsgenauigkeit. Es verallgemeinert am besten.
  Das tiefere Netz $A$ ist trotz höherer Trainingsgenauigkeit schlechter auf
  den Testdaten.
\end{enumerate}
}

% -------------------------------------------------------------------
% 3) Rechenaufwand einer Dense-Schicht
% -------------------------------------------------------------------
\aufgabe{
Die Ausgabe einer Dense-Schicht entsteht durch die Matrixmultiplikation
$\text{input}\cdot W$ mit
$\text{input}\in\mathbb{R}^{1\times n_{\text{in}}}$ und
$W\in\mathbb{R}^{n_{\text{in}}\times n_{\text{out}}}$. Jeder Ausgabewert
benötigt $n_{\text{in}}$ Multiplikationen (die Additionen und der Bias
werden hier vernachlässigt).

\begin{enumerate}[label=\alph*)]
  \item Wie viele Multiplikationen fallen für \textbf{ein} Bild in der
        Schicht $784 \rightarrow 512$ an?
  \item Berechnen Sie die Multiplikationen pro Bild für alle drei
        versteckten Schichten des tiefen Netzes
        $784 \rightarrow 512 \rightarrow 256 \rightarrow 128 \rightarrow 10$
        (inkl. Ausgabeschicht).
  \item Wie viele Multiplikationen fallen an, wenn ein Batch aus $600$
        Bildern durch das gesamte Netz läuft?
\end{enumerate}
}

\loesung{
\begin{enumerate}[label=\alph*)]
  \item Jede der $512$ Ausgaben braucht $784$ Multiplikationen:
  \[
    784\cdot 512 = 401\,408 .
  \]

  \item Pro Schicht $n_{\text{in}}\cdot n_{\text{out}}$:
  \[
    784\cdot 512 = 401\,408 ,
  \]
  \[
    512\cdot 256 = 131\,072 ,
  \]
  \[
    256\cdot 128 = 32\,768 ,
  \]
  \[
    128\cdot 10 = 1\,280 .
  \]
  Summe pro Bild:
  \[
    401\,408 + 131\,072 + 32\,768 + 1\,280 = 566\,528 .
  \]

  \item Für einen Batch aus $600$ Bildern:
  \[
    600\cdot 566\,528 = 339\,916\,800 \approx 3{,}4\cdot 10^{8}
    \text{ Multiplikationen} .
  \]
  Diese Operationen sind unabhängig und lassen sich auf der GPU
  \textbf{parallel} berechnen, was den Geschwindigkeitsvorteil erklärt.
\end{enumerate}
}

% -------------------------------------------------------------------
% 4) Parallelisierung: CPU gegen GPU
% -------------------------------------------------------------------
\aufgabe{
Eine Matrixmultiplikation erfordere $10^{9}$ elementare Operationen. Eine
CPU mit $8$ parallelen Recheneinheiten schafft $10^{9}$ Operationen pro
Einheit und Sekunde. Eine GPU besitzt $4\,096$ Recheneinheiten mit
jeweils $0{,}5\cdot 10^{9}$ Operationen pro Sekunde.

\begin{enumerate}[label=\alph*)]
  \item Berechnen Sie den Gesamtdurchsatz (Operationen pro Sekunde) von CPU
        und GPU.
  \item Wie lange braucht jede Hardware für die Matrixmultiplikation?
  \item Um welchen Faktor ist die GPU schneller?
\end{enumerate}
}

\loesung{
\begin{enumerate}[label=\alph*)]
  \item Durchsatz ist Anzahl Einheiten mal Leistung je Einheit:
  \[
    \text{CPU}: 8\cdot 10^{9} = 8\cdot 10^{9} \text{ Op/s} ,
  \]
  \[
    \text{GPU}: 4\,096\cdot 0{,}5\cdot 10^{9} = 2\,048\cdot 10^{9}
    = 2{,}048\cdot 10^{12} \text{ Op/s} .
  \]

  \item Zeit ist Operationen geteilt durch Durchsatz:
  \[
    t_{\text{CPU}} = \frac{10^{9}}{8\cdot 10^{9}} = 0{,}125 \text{ s} ,
  \]
  \[
    t_{\text{GPU}} = \frac{10^{9}}{2{,}048\cdot 10^{12}}
    \approx 4{,}88\cdot 10^{-4} \text{ s} .
  \]

  \item Beschleunigungsfaktor:
  \[
    \frac{t_{\text{CPU}}}{t_{\text{GPU}}}
    = \frac{0{,}125}{4{,}88\cdot 10^{-4}} \approx 256 .
  \]
  Der Faktor entspricht genau dem Verhältnis der Gesamtdurchsätze
  ($2\,048\cdot 10^{9} / (8\cdot 10^{9}) = 256$). Die GPU gewinnt durch ihre
  hohe Parallelität, obwohl jede einzelne Einheit langsamer rechnet.
\end{enumerate}
}

% -------------------------------------------------------------------
% 5) Wachstum des Rechenaufwands mit der Netztiefe
% -------------------------------------------------------------------
\aufgabe{
Ein Netz habe eine Eingabe der Größe $n$ und $L$ versteckte Schichten mit
je $n$ Neuronen (der Einfachheit halber gleich breit). Der Rechenaufwand
einer Schicht ist $n^2$ Multiplikationen.

\begin{enumerate}[label=\alph*)]
  \item Geben Sie den Gesamtaufwand des Netzes in Abhängigkeit von $n$ und
        $L$ an.
  \item Berechnen Sie den Aufwand für $n = 500$ und $L = 3$ sowie für
        $L = 6$.
  \item Wie verändert sich der Aufwand, wenn man bei fester Tiefe $L$ die
        Breite $n$ verdoppelt?
\end{enumerate}
}

\loesung{
\begin{enumerate}[label=\alph*)]
  \item Jede der $L$ Schichten kostet $n^2$, also
  \[
    \text{Aufwand} = L\cdot n^2 .
  \]

  \item Für $n = 500$ ist $n^2 = 250\,000$:
  \[
    L = 3: \quad 3\cdot 250\,000 = 750\,000 ,
  \]
  \[
    L = 6: \quad 6\cdot 250\,000 = 1\,500\,000 .
  \]
  Doppelte Tiefe bedeutet hier doppelten Aufwand (linear in $L$).

  \item Verdoppelt man $n$, so wächst $n^2$ auf $(2n)^2 = 4n^2$. Der Aufwand
  \textbf{vervierfacht} sich also. Die Breite geht \textbf{quadratisch},
  die Tiefe nur \textbf{linear} in den Aufwand ein.
\end{enumerate}
}


\subsection{CNNs}


% ===================================================================
% Aufgaben zu Convolutional Neural Networks (CNN)
% Themen: Graustufenumrechnung, 2D-Faltung, Ausgabegroesse,
%         Parameterzahl von Conv-Layern, MaxPooling,
%         rezeptives Feld, separierbare Filter, Dropout
% ===================================================================

% -------------------------------------------------------------------
% 1) RGB zu Graustufen
% -------------------------------------------------------------------
\aufgabe{
Die Umrechnung eines RGB-Pixels in einen Graustufenwert nutzt die
gewichtete Summe
\[
  I_{\text{grau}} = 0{,}299\,R + 0{,}587\,G + 0{,}114\,B .
\]
Berechnen Sie den Grauwert (auf ganze Zahlen gerundet) für:

\begin{enumerate}[label=\alph*)]
  \item reines Rot $(R,G,B) = (255, 0, 0)$,
  \item reines Grün $(0, 255, 0)$,
  \item einen Mischton $(100, 150, 200)$.
\end{enumerate}
}

\loesung{
\begin{enumerate}[label=\alph*)]
  \item
  \[
    I_{\text{grau}} = 0{,}299\cdot 255 = 76{,}2 \approx 76 .
  \]

  \item
  \[
    I_{\text{grau}} = 0{,}587\cdot 255 = 149{,}7 \approx 150 .
  \]
  Grün erscheint dem Auge deutlich heller als Rot, obwohl beide dieselbe
  Intensität $255$ haben.

  \item
  \[
    I_{\text{grau}} = 0{,}299\cdot 100 + 0{,}587\cdot 150 + 0{,}114\cdot 200
  \]
  \[
    = 29{,}9 + 88{,}05 + 22{,}8 = 140{,}75 \approx 141 .
  \]
\end{enumerate}
}

% -------------------------------------------------------------------
% 2) 2D-Faltung mit Boxfilter (Valid)
% -------------------------------------------------------------------
\aufgabe{
Gegeben sei der $4\times 4$-Bildausschnitt
\[
  I =
  \begin{pmatrix}
    2 & 4 & 6 & 8\\
    0 & 2 & 4 & 6\\
    8 & 6 & 4 & 2\\
    0 & 0 & 2 & 4
  \end{pmatrix}
\]
und der $3\times 3$-Mittelwertfilter
$k = \tfrac{1}{9}\begin{pmatrix} 1&1&1\\1&1&1\\1&1&1 \end{pmatrix}$. Es wird
mit \textbf{Stride 1} und \textbf{Valid}-Faltung (ohne Padding) gearbeitet.

\begin{enumerate}[label=\alph*)]
  \item Welche Größe hat das Ausgabebild?
  \item Berechnen Sie alle Werte des Ausgabebildes (auf zwei Nachkommastellen).
\end{enumerate}
}

\loesung{
\begin{enumerate}[label=\alph*)]
  \item Ausgabegröße $W - F + 1 = 4 - 3 + 1 = 2$, also ein $2\times 2$-Bild.

  \item Da der Filter symmetrisch ist, entspricht die Faltung dem Mittelwert
  jedes $3\times 3$-Fensters. Oben links (Zeilen 1--3, Spalten 1--3):
  \[
    \frac{2+4+6 + 0+2+4 + 8+6+4}{9} = \frac{36}{9} = 4{,}00 .
  \]
  Oben rechts (Zeilen 1--3, Spalten 2--4):
  \[
    \frac{4+6+8 + 2+4+6 + 6+4+2}{9} = \frac{42}{9} \approx 4{,}67 .
  \]
  Unten links (Zeilen 2--4, Spalten 1--3):
  \[
    \frac{0+2+4 + 8+6+4 + 0+0+2}{9} = \frac{26}{9} \approx 2{,}89 .
  \]
  Unten rechts (Zeilen 2--4, Spalten 2--4):
  \[
    \frac{2+4+6 + 6+4+2 + 0+2+4}{9} = \frac{30}{9} \approx 3{,}33 .
  \]
  \[
    O \approx
    \begin{pmatrix} 4{,}00 & 4{,}67\\ 2{,}89 & 3{,}33 \end{pmatrix} .
  \]
\end{enumerate}
}

% -------------------------------------------------------------------
% 3) Faltung mit Sobel-Filter
% -------------------------------------------------------------------
\aufgabe{
Der Sobel-Filter für horizontale Ableitungen ist
\[
  K_x =
  \begin{pmatrix} -1 & 0 & 1\\ -2 & 0 & 2\\ -1 & 0 & 1 \end{pmatrix} .
\]
Er wird als \textbf{Kreuzkorrelation} (ohne Kernel-Spiegelung) auf den
Bildausschnitt
\[
  A =
  \begin{pmatrix} 10 & 10 & 50\\ 10 & 10 & 50\\ 10 & 10 & 50 \end{pmatrix}
\]
angewendet.

\begin{enumerate}[label=\alph*)]
  \item Berechnen Sie den Ausgabewert an der Fenstermitte.
  \item Berechnen Sie denselben Wert für einen homogenen Ausschnitt, bei dem
        alle Einträge $= 10$ sind. Was zeigt das Ergebnis?
\end{enumerate}
}

\loesung{
\begin{enumerate}[label=\alph*)]
  \item Elementweise Multiplikation und Summe:
  \[
    (-1)\cdot 10 + 0\cdot 10 + 1\cdot 50
  \]
  \[
    + (-2)\cdot 10 + 0\cdot 10 + 2\cdot 50
  \]
  \[
    + (-1)\cdot 10 + 0\cdot 10 + 1\cdot 50 .
  \]
  Zeilenweise: $(-10 + 50) + (-20 + 100) + (-10 + 50) = 40 + 80 + 40 = 160$.
  Der große positive Wert zeigt eine \textbf{senkrechte Kante} (Sprung von
  $10$ auf $50$ nach rechts).

  \item Bei konstantem Ausschnitt heben sich die positiven und negativen
  Gewichte auf:
  \[
    (-1-2-1 + 1+2+1)\cdot 10 = 0 .
  \]
  In homogenen Regionen liefert der Sobel-Filter also $0$, passend zur
  Gewichtssumme $0$. Er reagiert nur auf Intensitätsänderungen.
\end{enumerate}
}

% -------------------------------------------------------------------
% 4) Ausgabegroesse mit Padding und Stride
% -------------------------------------------------------------------
\aufgabe{
Die Ausgabegröße einer Faltungsschicht ist
\[
  W_{\text{out}} = \left\lfloor \frac{W - F + 2P}{S} \right\rfloor + 1 ,
\]
mit Eingabebreite $W$, Kernelgröße $F$, Padding $P$ und Stride $S$.
Berechnen Sie $W_{\text{out}}$ für:

\begin{enumerate}[label=\alph*)]
  \item $W = 28$, $F = 3$, $P = 0$, $S = 1$ (Valid),
  \item $W = 28$, $F = 3$, $P = 1$, $S = 1$ (Same),
  \item $W = 32$, $F = 5$, $P = 0$, $S = 2$,
  \item Welches Padding $P$ ist für einen $5\times 5$-Kernel nötig, damit
        bei $S = 1$ die Größe erhalten bleibt (Same)?
\end{enumerate}
}

\loesung{
\begin{enumerate}[label=\alph*)]
  \item
  \[
    W_{\text{out}} = \left\lfloor \frac{28 - 3 + 0}{1} \right\rfloor + 1
    = 25 + 1 = 26 .
  \]

  \item
  \[
    W_{\text{out}} = \left\lfloor \frac{28 - 3 + 2}{1} \right\rfloor + 1
    = 27 + 1 = 28 .
  \]
  Das Padding $P = 1$ erhält bei einem $3\times 3$-Kernel die Größe.

  \item
  \[
    W_{\text{out}} = \left\lfloor \frac{32 - 5 + 0}{2} \right\rfloor + 1
    = \lfloor 13{,}5 \rfloor + 1 = 13 + 1 = 14 .
  \]

  \item Für Same-Padding gilt $P = \tfrac{F-1}{2} = \tfrac{5-1}{2} = 2$.
\end{enumerate}
}

% -------------------------------------------------------------------
% 5) Parameterzahl einer Conv2D-Schicht
% -------------------------------------------------------------------
\aufgabe{
Eine Conv2D-Schicht mit $D$ Kerneln der Größe $k\times k$ auf einem Eingang
mit $C$ Kanälen hat
\[
  \#\text{Parameter} = D\cdot(k^2\cdot C + 1)
\]
(der Summand $+1$ ist der Bias je Kernel).

\begin{enumerate}[label=\alph*)]
  \item Berechnen Sie die Parameterzahl für $k = 3$, $C = 1$, $D = 64$
        (erste Schicht des MNIST-CNN).
  \item Berechnen Sie sie für $k = 3$, $C = 32$, $D = 64$ (eine tiefere
        Schicht).
  \item Vergleichen Sie mit einer Dense-Schicht, die dieselben $28\cdot 28$
        Eingänge auf $64$ Neuronen abbildet. Warum sind CNNs
        parametersparsamer?
\end{enumerate}
}

\loesung{
\begin{enumerate}[label=\alph*)]
  \item
  \[
    64\cdot(3^2\cdot 1 + 1) = 64\cdot(9 + 1) = 64\cdot 10 = 640 .
  \]

  \item
  \[
    64\cdot(3^2\cdot 32 + 1) = 64\cdot(288 + 1) = 64\cdot 289 = 18\,496 .
  \]

  \item Dense-Schicht $784 \rightarrow 64$:
  \[
    784\cdot 64 + 64 = 50\,176 + 64 = 50\,240 \text{ Parameter} .
  \]
  Die Conv-Schicht braucht nur $640$ Parameter, weil derselbe kleine Kernel
  über \textbf{alle} Bildpositionen wiederverwendet wird
  (\textbf{Gewichtsteilung}). Bei Dense hat jede Position eigene Gewichte.
\end{enumerate}
}

% -------------------------------------------------------------------
% 6) Parameter des einfachen MNIST-CNN
% -------------------------------------------------------------------
\aufgabe{
Ein CNN für MNIST hat den Aufbau
\[
  \text{Input } (28,28,1) \;\rightarrow\;
  \text{Conv2D}(64, 3\times 3) \;\rightarrow\;
  \text{Flatten} \;\rightarrow\;
  \text{Dense}(128) \;\rightarrow\;
  \text{Dense}(10) .
\]
Die Conv-Schicht nutzt Valid-Faltung (kein Padding, Stride 1).

\begin{enumerate}[label=\alph*)]
  \item Bestimmen Sie die Ausgabeform der Conv2D-Schicht und die Anzahl der
        Werte nach dem Flatten.
  \item Berechnen Sie die Parameterzahl jeder Schicht.
  \item In welcher Schicht liegen die meisten Parameter, und warum?
\end{enumerate}
}

\loesung{
\begin{enumerate}[label=\alph*)]
  \item Conv verkleinert $28$ auf $28 - 3 + 1 = 26$, mit $64$ Filtern:
  Ausgabeform $(26, 26, 64)$. Flatten ergibt
  \[
    26\cdot 26\cdot 64 = 43\,264 \text{ Werte} .
  \]

  \item Parameter:
  \[
    \text{Conv2D}: 64\cdot(3^2\cdot 1 + 1) = 640 ,
  \]
  \[
    \text{Dense}(128): 43\,264\cdot 128 + 128 = 5\,537\,920 ,
  \]
  \[
    \text{Dense}(10): 128\cdot 10 + 10 = 1\,290 .
  \]
  Flatten hat keine Parameter. Gesamt:
  \[
    640 + 5\,537\,920 + 1\,290 = 5\,539\,850 .
  \]

  \item Die erste Dense-Schicht enthält über $99\,\%$ aller Parameter. Nach
  dem Flatten sind $43\,264$ Eingänge mit $128$ Neuronen vollverbunden. Ein
  \textbf{Pooling} vor dem Flatten würde diese Zahl drastisch senken.
\end{enumerate}
}

% -------------------------------------------------------------------
% 7) MaxPooling
% -------------------------------------------------------------------
\aufgabe{
MaxPooling mit Fenster $2\times 2$ und Stride $2$ wählt aus jedem Fenster den
größten Wert. Gegeben sei die Feature Map
\[
  F =
  \begin{pmatrix}
    1 & 3 & 2 & 4\\
    5 & 6 & 1 & 0\\
    2 & 0 & 7 & 8\\
    1 & 2 & 3 & 4
  \end{pmatrix}
\]

\begin{enumerate}[label=\alph*)]
  \item Berechnen Sie die gepoolte Ausgabe.
  \item Um welchen Faktor sinkt die Anzahl der Werte?
  \item Warum hat eine MaxPooling-Schicht keine trainierbaren Parameter?
\end{enumerate}
}

\loesung{
\begin{enumerate}[label=\alph*)]
  \item Vier nicht überlappende Fenster:
  \[
    \max(1,3,5,6) = 6, \qquad \max(2,4,1,0) = 4,
  \]
  \[
    \max(2,0,1,2) = 2, \qquad \max(7,8,3,4) = 8 .
  \]
  \[
    O = \begin{pmatrix} 6 & 4\\ 2 & 8 \end{pmatrix} .
  \]

  \item Aus $4\times 4 = 16$ Werten werden $2\times 2 = 4$, also eine
  Reduktion auf ein Viertel ($75\,\%$ weniger).

  \item MaxPooling wählt nur den größten Wert aus, es gibt keine Gewichte,
  die gelernt werden könnten. Die Operation ist fest vorgegeben.
\end{enumerate}
}

% -------------------------------------------------------------------
% 8) Rezeptives Feld und zwei 3x3 statt einer 5x5
% -------------------------------------------------------------------
\aufgabe{
Bei $3\times 3$-Faltungen mit Stride $1$ vergrößert jede weitere Schicht das
rezeptive Feld um $2$ Pixel je Kante (die erste Schicht hat $3\times 3$).

\begin{enumerate}[label=\alph*)]
  \item Wie groß ist das rezeptive Feld nach zwei bzw. drei gestapelten
        $3\times 3$-Faltungen?
  \item Vergleichen Sie die Gewichtszahl von zwei $3\times 3$-Faltungen mit
        einer $5\times 5$-Faltung bei jeweils $C$ Ein- und Ausgabekanälen
        (ohne Bias).
  \item Warum ist der Stapel aus zwei $3\times 3$-Schichten trotz gleichem
        rezeptiven Feld ausdrucksstärker?
\end{enumerate}
}

\loesung{
\begin{enumerate}[label=\alph*)]
  \item Erste Schicht $3\times 3$, jede weitere $+2$:
  \[
    \text{zwei Schichten}: 5\times 5, \qquad
    \text{drei Schichten}: 7\times 7 .
  \]

  \item Gewichte pro Schicht $k^2\cdot C^2$:
  \[
    \text{eine } 5\times 5: \; 25\,C^2, \qquad
    \text{zwei } 3\times 3: \; 2\cdot 9\,C^2 = 18\,C^2 .
  \]
  Der Stapel spart $\tfrac{25 - 18}{25} = 28\,\%$ der Gewichte.

  \item Zwischen den beiden $3\times 3$-Faltungen liegt eine
  \textbf{ReLU}-Nichtlinearität. Ohne sie ließen sich beide Schichten zu
  einer einzigen linearen Faltung zusammenfassen. Die zusätzliche
  Nichtlinearität erlaubt komplexere Funktionen bei gleichem rezeptiven Feld
  und weniger Parametern.
\end{enumerate}
}

% -------------------------------------------------------------------
% 9) Separierbare Filter
% -------------------------------------------------------------------
\aufgabe{
Ein Kernel heißt separierbar, wenn er sich als äußeres Produkt
$K = k_y\,k_x^{\mathsf{T}}$ zweier Vektoren schreiben lässt. Betrachtet wird
der Sobel-Kernel
\[
  K =
  \begin{pmatrix} -1 & 0 & 1\\ -2 & 0 & 2\\ -1 & 0 & 1 \end{pmatrix},
  \qquad
  k_y = \begin{pmatrix} 1 \\ 2 \\ 1 \end{pmatrix},
  \qquad
  k_x = \begin{pmatrix} -1 \\ 0 \\ 1 \end{pmatrix} .
\]

\begin{enumerate}[label=\alph*)]
  \item Bestätigen Sie durch Berechnung von $k_y\,k_x^{\mathsf{T}}$, dass
        $K$ separierbar ist.
  \item Ein $k\times k$-Filter kostet pro Pixel $k^2$ Multiplikationen, die
        separierte Variante nur $2k$. Berechnen Sie beide Werte für
        $k = 3$ und $k = 7$.
\end{enumerate}
}

\loesung{
\begin{enumerate}[label=\alph*)]
  \item Äußeres Produkt (Spaltenvektor mal Zeilenvektor):
  \[
    k_y\,k_x^{\mathsf{T}} =
    \begin{pmatrix} 1 \\ 2 \\ 1 \end{pmatrix}
    \begin{pmatrix} -1 & 0 & 1 \end{pmatrix}
    =
    \begin{pmatrix}
      -1 & 0 & 1\\
      -2 & 0 & 2\\
      -1 & 0 & 1
    \end{pmatrix}
    = K .
  \]
  Die Darstellung stimmt, $K$ ist separierbar.

  \item Aufwand pro Pixel:
  \[
    k = 3: \quad k^2 = 9 \;\text{ gegen }\; 2k = 6 ,
  \]
  \[
    k = 7: \quad k^2 = 49 \;\text{ gegen }\; 2k = 14 .
  \]
  Der Vorteil wächst mit der Kernelgröße: Bei $k = 7$ spart die separierte
  Faltung fast drei Viertel der Multiplikationen.
\end{enumerate}
}

% -------------------------------------------------------------------
% 10) Dropout
% -------------------------------------------------------------------
\aufgabe{
Eine Dropout-Schicht mit Rate $p$ setzt jeden Aktivierungswert mit
Wahrscheinlichkeit $p$ auf $0$ und skaliert die überlebenden Werte im
Training mit dem Faktor $\tfrac{1}{1-p}$ (Inverted Dropout).

\begin{enumerate}[label=\alph*)]
  \item Berechnen Sie den Skalierungsfaktor für $p = 0{,}25$ und $p = 0{,}5$.
  \item Gegeben seien die Aktivierungen $(2{,}0;\ 4{,}0;\ 1{,}0;\ 3{,}0)$ und
        $p = 0{,}5$. Der erste und dritte Wert fallen aus. Geben Sie den
        Ausgabevektor an.
  \item Zeigen Sie, dass der Erwartungswert einer Aktivierung durch die
        Skalierung erhalten bleibt.
\end{enumerate}
}

\loesung{
\begin{enumerate}[label=\alph*)]
  \item Faktor $\tfrac{1}{1-p}$:
  \[
    p = 0{,}25: \; \frac{1}{0{,}75} \approx 1{,}33, \qquad
    p = 0{,}5: \; \frac{1}{0{,}5} = 2 .
  \]

  \item Ausgefallene Werte werden $0$, die überlebenden verdoppelt:
  \[
    (0;\ 8{,}0;\ 0;\ 6{,}0) .
  \]

  \item Sei $a$ eine Aktivierung. Mit Wahrscheinlichkeit $1-p$ überlebt sie
  und wird zu $\tfrac{a}{1-p}$, mit Wahrscheinlichkeit $p$ wird sie $0$:
  \[
    E[\text{Ausgabe}]
    = (1-p)\cdot\frac{a}{1-p} + p\cdot 0 = a .
  \]
  Der Erwartungswert entspricht also der ursprünglichen Aktivierung. Das ist
  der Grund für die Skalierung, so bleibt das Netz im Mittel gleich stark
  angeregt wie ohne Dropout.
\end{enumerate}
}